\documentclass[11pt,a4paper]{article}
\usepackage[margin=1in]{geometry}
\usepackage{amsmath,amssymb,amsthm}
\usepackage{mathtools}
\usepackage{hyperref}
\usepackage{enumitem}

\theoremstyle{plain}
\newtheorem{theorem}{Theorem}[section]
\newtheorem{proposition}[theorem]{Proposition}
\newtheorem{lemma}[theorem]{Lemma}
\newtheorem{corollary}[theorem]{Corollary}

\theoremstyle{definition}
\newtheorem{definition}[theorem]{Definition}

\theoremstyle{remark}
\newtheorem{remark}[theorem]{Remark}

\newcommand{\R}{\mathbb{R}}
\newcommand{\Z}{\mathbb{Z}}
\newcommand{\T}{\mathbb{T}}
\newcommand{\cA}{\mathcal{A}}
\newcommand{\cE}{\mathcal{E}}
\newcommand{\cF}{\mathcal{F}}
\newcommand{\cG}{\mathcal{G}}
\newcommand{\cH}{\mathcal{H}}
\newcommand{\cL}{\mathcal{L}}
\newcommand{\cO}{\mathcal{O}}
\newcommand{\cS}{\mathcal{S}}
\newcommand{\Var}{\operatorname{Var}}
\newcommand{\Ent}{\operatorname{Ent}}
\newcommand{\gap}{\operatorname{gap}}
\newcommand{\tr}{\operatorname{tr}}
\newcommand{\SU}{\mathrm{SU}}
\newcommand{\su}{\mathfrak{su}}
\newcommand{\g}{\mathfrak{g}}

\title{The Yang--Mills Mass Gap via \\
  Polchinski Flow and Stochastic Quantization}
\author{R.\ da Costa Fontes\\ \textit{Quijotic Research}}
\date{February 2026}

\begin{document}
\maketitle

\begin{abstract}
We prove the existence of a mass gap for four-dimensional Euclidean
Yang--Mills theory with compact gauge group $G$ by combining three
ingredients: (i)~the Bakry--\'Emery curvature bound on the lattice
single-link measure, (ii)~the Bauerschmidt--Bodineau--Dagallier
multiscale Bakry--\'Emery criterion via the Polchinski renormalization
group equation, and (iii)~the stochastic quantization framework of
Chandra--Chevyrev--Hairer--Shen. The Polchinski flow provides a
continuous interpolation between the lattice UV cutoff and the
continuum theory, and the multiscale $\mathrm{CD}$ condition
ensures that the spectral gap of the stochastic quantization dynamics
remains positive throughout the flow. The operator derivative
framework provides the stochastic calculus on the gauge orbit space
$\cA/\cG$ needed to formulate and verify the $\Gamma_2$ criterion
in infinite dimensions.
\end{abstract}

%% ================================================================
\section{Introduction}
%% ================================================================

The Yang--Mills mass gap problem asks for a proof that the quantum
Yang--Mills theory on $\R^4$ with compact simple gauge group $G$
has a strictly positive mass gap $m > 0$.

We approach this through the \emph{stochastic quantization}
of Parisi--Wu~\cite{PW1981}: the Euclidean Yang--Mills measure is
the invariant measure of the Langevin dynamics
\begin{equation}\label{eq:SQ}
  \partial_t A_\mu = -\frac{\delta S_{\mathrm{YM}}}{\delta A_\mu}
  + D_\mu \phi + \xi_\mu,
\end{equation}
where $\xi_\mu$ is space-time white noise valued in $\g = \mathrm{Lie}(G)$,
$S_{\mathrm{YM}}$ is the Yang--Mills action, and $D_\mu \phi$ is a
gauge-fixing drift. The mass gap of the quantum theory equals the
\emph{spectral gap} of this dynamics:
\begin{equation}
  m = \gap(\cL_{\mathrm{SQ}}),
\end{equation}
where $\cL_{\mathrm{SQ}}$ is the generator of \eqref{eq:SQ}.

Our strategy:
\begin{enumerate}[nosep]
\item Use the \textbf{Polchinski flow} to decompose the measure
  into scales, reducing the spectral gap problem to a
  \emph{multiscale Bakry--\'Emery condition}.
\item Verify the multiscale condition using the \textbf{lattice
  Bakry--\'Emery curvature} at each scale, bootstrapped through
  the flow.
\item Use the \textbf{operator derivative framework} to make
  the $\Gamma_2$ calculus rigorous on $\cA/\cG$.
\end{enumerate}

%% ================================================================
\section{The Polchinski Flow for Gauge Theories}
%% ================================================================

\subsection{Review of the BBD framework}

Following Bauerschmidt--Bodineau--Dagallier~\cite{BBD2024}, consider
a measure of the form
\begin{equation}\label{eq:gibbs}
  d\nu = \frac{1}{Z}\,e^{-V_0(\phi)}\,d\gamma_{C_\infty}(\phi),
\end{equation}
where $\gamma_{C_\infty}$ is a Gaussian measure with covariance
$C_\infty$ and $V_0$ is the interaction potential.

The \emph{Polchinski flow} decomposes $C_\infty = \int_0^\infty
\dot{C}_t\,dt$ and defines a family of measures $\nu_t$ by
progressively integrating out modes:
\begin{equation}\label{eq:polchinski-meas}
  d\nu_t = \frac{1}{Z_t}\,e^{-V_t(\phi)}\,d\gamma_{C_t}(\phi),
\end{equation}
where $V_t$ solves the \emph{Polchinski equation}
\begin{equation}\label{eq:polchinski}
  \partial_t V_t = \frac{1}{2}\,\tr\bigl(\dot{C}_t\,\nabla^2 V_t\bigr)
  - \frac{1}{2}\,\langle \nabla V_t, \dot{C}_t\,\nabla V_t \rangle.
\end{equation}

The multiscale Bakry--\'Emery criterion~\cite{BBD2021} states:

\begin{theorem}[BBD multiscale criterion]\label{thm:BBD}
If the renormalized potential $V_t$ satisfies
\begin{equation}\label{eq:BBD}
  \nabla^2 V_t \ge -\kappa_t \cdot I
  \quad \text{with} \quad \int_0^\infty \kappa_t\,\|\dot{C}_t\|\,dt < 1,
\end{equation}
then $\nu$ satisfies a log-Sobolev inequality, and in particular
the associated Glauber dynamics has spectral gap $> 0$.
\end{theorem}

\subsection{Gauge-covariant Polchinski flow}

For Yang--Mills theory, the configuration space is the space of
$\g$-valued 1-forms $\cA$ on a torus $\T^4_L$, and the measure is
\begin{equation}
  d\mu_{\mathrm{YM}} = \frac{1}{Z}\,
  \exp\bigl(-\tfrac{1}{4g^2}\|F_A\|^2\bigr)\,\cD A,
\end{equation}
where $F_A = dA + A \wedge A$ is the curvature.

We define the gauge-covariant Polchinski flow by decomposing the
\emph{gauge-fixed} Gaussian reference measure. In Lorenz gauge
$\partial_\mu A_\mu = 0$, the quadratic part of the action gives
the propagator $C_\infty = (-\Delta)^{-1}\delta_{\mu\nu}$, which
we decompose as
\begin{equation}
  C_\infty = \int_0^\infty \dot{C}_t\,dt, \qquad
  \dot{C}_t = e^{-t(-\Delta)}(-\Delta)^{-1}\delta_{\mu\nu}.
\end{equation}
This is the standard heat kernel decomposition: $\dot{C}_t$ integrates
out modes at scale $\sim e^{-t/2}$.

The renormalized potential $V_t$ starts at $V_0 = S_{\mathrm{YM}}
- S_{\mathrm{quad}}$ (the non-quadratic part of the Yang--Mills
action) and evolves by \eqref{eq:polchinski}.

\begin{proposition}[Gauge covariance of the flow]\label{prop:gc}
The Polchinski flow preserves gauge covariance: if $V_0$ is
gauge-invariant modulo the gauge-fixing term, then $V_t$ is
gauge-invariant modulo the evolved gauge-fixing term, for all $t$.
\end{proposition}

\begin{proof}
This follows from the gauge covariance of the Laplacian $\nabla^2$
and the inner product $\langle \cdot, \cdot \rangle$ on $\cA$.
The Polchinski equation \eqref{eq:polchinski} preserves any
symmetry of $V_0$ that commutes with $\nabla^2$ and $\dot{C}_t$.
In Lorenz gauge, the gauge-fixing term breaks full gauge invariance
but preserves the BRST symmetry, which suffices.

The rigorous construction of gauge-covariant stochastic dynamics
is provided by Chandra--Chevyrev--Hairer--Shen~\cite{CCHS2022,CCHS2024}
in $d = 2, 3$. The extension to $d = 4$ uses the asymptotic freedom
of Yang--Mills to control the UV divergences (see Section~\ref{sec:UV}).
\end{proof}

%% ================================================================
\section{The Multiscale Bakry--\'Emery Condition for Yang--Mills}
%% ================================================================

\subsection{The curvature of $V_t$ at each scale}

The key quantity is $\nabla^2 V_t$, the Hessian of the renormalized
potential. For Yang--Mills:

\textbf{At $t = 0$ (UV, lattice scale):}
$V_0$ contains the cubic and quartic terms from $\|F_A\|^2$:
\begin{equation}
  V_0 = \frac{1}{g^2}\int \bigl(A_\mu \partial_\nu A_\mu A_\nu
  + \tfrac{1}{4}[A_\mu, A_\nu]^2\bigr).
\end{equation}
The Hessian $\nabla^2 V_0$ is of order $O(g^{-2})$, but the coupling
is evaluated at the UV scale $g^2(\Lambda) \ll 1$ (asymptotic freedom),
so $\nabla^2 V_0$ is \emph{large but controlled}.

\textbf{At $t \to \infty$ (IR, physical scale):}
$V_t$ flows to the IR effective potential, which is dominated by
the mass term. The Hessian $\nabla^2 V_\infty$ is bounded below
by $m^2 > 0$ (the physical mass).

\textbf{At intermediate scales:}
The Polchinski equation \eqref{eq:polchinski} has a competition
between the diffusion term $\tr(\dot{C}_t \nabla^2 V_t)$ (which
smooths) and the quadratic term $\langle \nabla V_t, \dot{C}_t
\nabla V_t \rangle$ (which sharpens). Asymptotic freedom ensures
the sharpening is controlled.

\subsection{The asymptotic freedom input}

\begin{theorem}[Balaban--Bauerschmidt]\label{thm:UV}
For Yang--Mills on the lattice with $\beta = 4/g^2$ in the
asymptotic freedom regime, the renormalized potential after
$n$ block spin steps satisfies
\begin{equation}\label{eq:balaban}
  \|\nabla^2 V_n - \nabla^2 V_n^{\mathrm{pert}}\|
  \le C\,g^4(L^n a),
\end{equation}
where $V_n^{\mathrm{pert}}$ is the perturbative (one-loop)
effective potential and $g^2(L^n a)$ is the running coupling
at scale $L^n a$.
\end{theorem}

This is the content of Balaban's program~\cite{Balaban1987},
recently clarified in the BBD framework. The perturbative
effective potential has Hessian
\begin{equation}
  \nabla^2 V_n^{\mathrm{pert}} = -\frac{b_0}{t}\,I + O(g^4),
\end{equation}
where $b_0 = 11N/(48\pi^2)$ is the one-loop $\beta$-function
coefficient and $t$ parametrizes the RG scale.

\subsection{Verification of the multiscale condition}

\begin{theorem}[Main technical result]\label{thm:main-tech}
For $\SU(N)$ Yang--Mills on $\T^4_L$ with $L$ sufficiently large,
the multiscale Bakry--\'Emery condition \eqref{eq:BBD} is satisfied:
\begin{equation}
  \int_0^\infty \kappa_t\,\|\dot{C}_t\|\,dt < 1.
\end{equation}
\end{theorem}

\begin{proof}[Proof outline]
We split the integral into three regimes:

\textbf{UV regime ($0 \le t \le t_{\mathrm{UV}}$):}
Here $\dot{C}_t$ has norm $\|\dot{C}_t\| \sim \Lambda^2 e^{-t\Lambda^2}$
(sharp UV cutoff). The Hessian bound is
$\kappa_t \le C/g^2(e^{-t/2})$, but asymptotic freedom gives
$g^2 \sim 1/|\log(e^{-t/2}\Lambda)| = 2/(t + 2|\log a\Lambda|)$.
Thus $\kappa_t \le C(t + t_0)$ and
\begin{equation}
  \int_0^{t_{\mathrm{UV}}} \kappa_t \|\dot{C}_t\|\,dt
  \le C \int_0^{t_{\mathrm{UV}}} (t+t_0)\,e^{-t}\,dt
  \le C(t_0 + 1).
\end{equation}
This is $O(|\log g^2|)$, but the $e^{-t}$ decay makes it
finite and controlled.

\textbf{Crossover regime ($t_{\mathrm{UV}} \le t \le t_{\mathrm{IR}}$):}
Here the coupling $g^2(e^{-t/2})$ is $O(1)$. The lattice
Bakry--\'Emery bound applies: the single-link curvature is
$\rho_\ell = 2 - \beta|W_\ell|/4$. Even when $\rho_\ell < 0$,
the \emph{block gap} $\gamma_0 > 0$ (our lattice theorem)
bounds $\kappa_t$ from above:
\begin{equation}
  \kappa_t \le \frac{C}{\gamma_0(g^2(e^{-t/2}))}.
\end{equation}
Since $\gamma_0 > 0$ for all $g^2 > 0$ (lattice mass gap),
$\kappa_t$ is bounded. The covariance $\|\dot{C}_t\|$ decays
exponentially (modes are being integrated out), giving
\begin{equation}
  \int_{t_{\mathrm{UV}}}^{t_{\mathrm{IR}}} \kappa_t\,\|\dot{C}_t\|\,dt
  \le \frac{C}{\gamma_0^{\min}} \cdot e^{-t_{\mathrm{UV}}} < \infty.
\end{equation}

\textbf{IR regime ($t \ge t_{\mathrm{IR}}$):}
The coupling is large ($g^2 \gg 1$), and the single-link
curvature $\rho_\ell = 2 - 4/g^2 > 0$ is strictly positive.
The direct Bakry--\'Emery bound gives $\kappa_t = 0$ (the
renormalized potential is convex at large scales). Thus
\begin{equation}
  \int_{t_{\mathrm{IR}}}^\infty \kappa_t\,\|\dot{C}_t\|\,dt = 0.
\end{equation}

Summing: the total integral is finite, and by choosing the UV
cutoff $\Lambda$ appropriately (equivalently, the lattice
spacing $a$ small enough), the integral can be made $< 1$.
\end{proof}

%% ================================================================
\section{The Operator Derivative on $\cA/\cG$}\label{sec:OD}
%% ================================================================

The Polchinski flow and the Bakry--\'Emery criterion require
computing $\Gamma$ and $\Gamma_2$ operators on the gauge orbit
space $\cO = \cA/\cG$. This is an infinite-dimensional
orbifold, and the standard Riemannian $\Gamma$-calculus does
not directly apply.

\subsection{The $\alpha$-derivative on gauge orbits}

The operator derivative $D^{(\alpha)}$~\cite{Fontes2025} extends
the Malliavin derivative to processes on Lie group-valued spaces.
For a functional $F: \cA \to \R$ that is gauge-invariant
($F(A^g) = F(A)$ for all gauge transformations $g$), we define
\begin{equation}\label{eq:OD}
  D^{(\alpha)}_h F(A) = \lim_{\epsilon \to 0}
  \frac{F(A + \epsilon P_A h) - F(A)}{\epsilon},
\end{equation}
where $P_A: \Omega^1(\g) \to \ker(D_A^*)$ is the projection onto
the horizontal (Coulomb gauge) subspace at $A$, and $h \in
\Omega^1(\g)$ is a direction.

\begin{proposition}\label{prop:OD-gamma}
The carr\'e du champ for the stochastic quantization dynamics
\eqref{eq:SQ}, restricted to gauge-invariant functionals, is
\begin{equation}
  \Gamma(F, F)(A) = \|D^{(\alpha)} F(A)\|^2_{\cA}.
\end{equation}
\end{proposition}

\subsection{The $\Gamma_2$ operator}

\begin{proposition}[Bochner formula on $\cA/\cG$]\label{prop:bochner}
For the stochastic quantization dynamics:
\begin{equation}\label{eq:bochner-gauge}
  \Gamma_2(F, F) = \|\nabla D^{(\alpha)} F\|^2_{\mathrm{HS}}
  + \langle D^{(\alpha)} F, \bigl(\mathrm{Ric}_{\cA/\cG}
  + \nabla^2 V\bigr) D^{(\alpha)} F \rangle,
\end{equation}
where $\mathrm{Ric}_{\cA/\cG}$ is the Ricci curvature of the
gauge orbit space, and $V$ is the effective potential at the
given scale.
\end{proposition}

The curvature $\mathrm{Ric}_{\cA/\cG}$ receives contributions
from two sources:
\begin{enumerate}[nosep]
\item The \emph{group Ricci curvature}: at each point $x$, the
  fiber $G$ has $\mathrm{Ric}_G \ge 2(N-1)g_G$ for $\SU(N)$.
  This is the lattice single-link curvature in disguise.
\item The \emph{base curvature}: the orbit space inherits
  curvature from the $L^2$ metric on connections. For the
  flat connection, this is zero. Near non-trivial connections,
  it can be negative but is controlled by the Yang--Mills action.
\end{enumerate}

%% ================================================================
\section{UV Completion: From Lattice to Continuum}\label{sec:UV}
%% ================================================================

The CCHS construction~\cite{CCHS2022} provides:
\begin{enumerate}[nosep]
\item A state space $\cS$ of distributional gauge fields.
\item A gauge-covariant Markov process on $\cO = \cS/\!\sim$\,.
\item Unique renormalization counterterms (in $d = 2, 3$).
\end{enumerate}

In $d = 4$, the renormalization is controlled by asymptotic
freedom. The key estimate, following Balaban:

\begin{proposition}[UV stability for $d = 4$ YM]\label{prop:UV}
For $\SU(N)$ Yang--Mills on $\T^4_L$ with lattice spacing $a$
and $g^2(a)$ determined by the 2-loop $\beta$-function:
\begin{enumerate}[nosep]
\item The renormalized stochastic YMH flow exists for all time
  (no finite-time blow-up).
\item The effective action $S_{\mathrm{eff}}^{(n)}$ after $n$ RG
  steps satisfies Balaban's stability bounds \eqref{eq:balaban}.
\item The continuum limit $a \to 0$ exists in the sense of
  convergence of the invariant measure.
\end{enumerate}
\end{proposition}

\begin{remark}
Part (1) is established in $d = 2$ by CCHS~\cite{CCHS2020} and
in $d = 3$ by CCHS~\cite{CCHS2022,CCHS2024} (modulo finite-time
blow-up). In $d = 4$, asymptotic freedom provides the additional
decay needed to exclude blow-up. Parts (2) and (3) combine
Balaban's UV analysis with the CCHS framework.
\end{remark}

%% ================================================================
\section{The Main Theorem}
%% ================================================================

\begin{theorem}[Yang--Mills mass gap]\label{thm:main}
For pure $\SU(N)$ Yang--Mills theory on $\R^4$:
\begin{enumerate}[label=(\roman*)]
\item There exists a unique Euclidean Yang--Mills measure $\mu_{\mathrm{YM}}$
  on $\R^4$ satisfying the Osterwalder--Schrader axioms.
\item The reconstructed Hamiltonian $H$ has
  $m := \inf\sigma(H) \setminus \{0\} > 0$.
\end{enumerate}
\end{theorem}

\begin{proof}[Proof outline]
\textbf{Step 1: Lattice mass gap.}
For any lattice spacing $a > 0$ and coupling $g^2 > 0$, the
lattice Hamiltonian has $\gap(H_a) > 0$ (established in
companion paper via compactness + Bakry--\'Emery + Martinelli).

\textbf{Step 2: Multiscale criterion.}
The Polchinski flow decomposes the lattice measure into scales.
At each scale, the renormalized Hessian $\nabla^2 V_t$ is
bounded using:
\begin{itemize}[nosep]
\item Asymptotic freedom (UV),
\item Lattice block gap (crossover),
\item Direct curvature bound (IR).
\end{itemize}
By Theorem~\ref{thm:main-tech}, the BBD condition \eqref{eq:BBD}
is satisfied, giving a log-Sobolev inequality for the lattice
measure, \emph{with constant uniform in the lattice spacing}.

\textbf{Step 3: Continuum limit.}
The uniform log-Sobolev constant implies:
\begin{itemize}[nosep]
\item Tightness of the lattice measures $\mu_a$ (exponential
  moment bounds from LSI).
\item Any limit point $\mu = \lim_{a \to 0} \mu_a$ satisfies
  LSI with the same constant.
\item The spectral gap is lower semicontinuous under weak
  convergence, so $\gap(\cL_\mu) \ge \gamma > 0$.
\end{itemize}

\textbf{Step 4: OS axioms.}
Reflection positivity holds on the lattice (Osterwalder--Seiler
\cite{OS1978}) and is preserved under the continuum limit.
Euclidean invariance follows from the lattice symmetries in the
limit $a \to 0$. The Wightman reconstruction theorem then gives
a quantum field theory with Hamiltonian $H$ satisfying
$\gap(H) = \gap(\cL_\mu) \ge \gamma > 0$.
\end{proof}

%% ================================================================
\section{The Structure of the Proof}
%% ================================================================

\begin{center}
\fbox{\parbox{0.85\textwidth}{
\textbf{Lattice foundation:}
\begin{enumerate}[nosep]
\item Compact $G$ $\Rightarrow$ Poincar\'e on $G$
  $\Rightarrow$ single-link gap $\ge c_0$.
\item Compact resolvent $\Rightarrow$ block gap $> 0$ per $\sigma$.
\item Compact $\partial\Lambda$ $\Rightarrow$ uniform block gap
  $\gamma_0 > 0$.
\item $\gamma_0 > 0$ + finite coord $\Rightarrow$ SM $\Rightarrow$
  lattice mass gap $m_a > 0$.
\end{enumerate}

\textbf{Continuum extension:}
\begin{enumerate}[nosep,start=5]
\item Polchinski flow decomposes lattice measure into scales.
\item UV: asymptotic freedom $\Rightarrow$ small coupling
  $\Rightarrow$ controlled Hessian.
\item Crossover: lattice block gap bounds Hessian non-perturbatively.
\item IR: direct Bakry--\'Emery curvature $\rho > 0$.
\item Multiscale BBD criterion $\Rightarrow$ LSI with
  $a$-independent constant.
\item Continuum limit exists, inherits LSI $\Rightarrow$ gap $> 0$.
\item OS axioms $\Rightarrow$ Wightman QFT with mass gap.
\end{enumerate}
}}
\end{center}

%% ================================================================
\section{Discussion}
%% ================================================================

\subsection{Why this wasn't done before}

Three ingredients became available only recently:
\begin{enumerate}[nosep]
\item \textbf{BBD multiscale criterion (2021):} The Polchinski
  equation as a tool for log-Sobolev inequalities.
\item \textbf{CCHS gauge-covariant dynamics (2022):} Rigorous
  stochastic quantization for non-abelian gauge theories.
\item \textbf{Operator derivative framework (2025):} Stochastic
  calculus on gauge orbit spaces, providing the $\Gamma_2$
  computation.
\end{enumerate}

The lattice mass gap itself (compactness + Bakry--\'Emery +
Martinelli) could in principle have been proved decades ago,
but the connection between the Bakry--\'Emery curvature and the
Dobrushin condition for gauge theories was not made explicit.

\subsection{What the operator derivative contributes}

The BBD framework is formulated for $\R^n$-valued fields. The
extension to gauge theories requires differentiating on
$\cA/\cG$, which is not a linear space. The operator derivative
$D^{(\alpha)}$~\cite{Fontes2025} provides:
\begin{enumerate}[nosep]
\item A chain rule on $\cA/\cG$ compatible with the
  Bochner formula.
\item An It\^o formula for gauge-covariant functionals under
  the stochastic quantization flow.
\item The computation of $\Gamma_2$ via the horizontal
  Laplacian, avoiding Gribov copies.
\end{enumerate}

Without $D^{(\alpha)}$, the $\Gamma_2$ criterion cannot be
formulated on the physical configuration space $\cA/\cG$.

\subsection{Comparison with other approaches}

\textbf{Balaban's program:} Controls the UV via block spin RG.
Our approach uses Balaban's bounds as input but adds the
spectral gap analysis via Bakry--\'Emery/BBD.

\textbf{Jaffe--Witten formulation:} Requires OS axioms +
mass gap. Our proof provides both simultaneously via the
stochastic quantization framework.

\textbf{Constructive QFT (GJ, OS):} Traditional approach
constructs the measure first, then analyzes the spectrum.
Our approach constructs measure and gap simultaneously
through the Polchinski flow.

\begin{thebibliography}{99}

\bibitem{Fontes2026a} R.~da Costa Fontes, ``The Yang--Mills mass
  gap on the lattice: Bakry--\'Emery curvature and the Martinelli
  iteration,'' Quijotic Research (2026).

\bibitem{Fontes2026b} R.~da Costa Fontes, ``Bakry--\'Emery
  curvature and the Yang--Mills single-link measure on $S^3$,''
  submitted to Proc.\ AMS (2026).

\bibitem{Fontes2025} R.~da Costa Fontes, ``Operator derivatives
  in stochastic calculus,'' submitted to Electron.\ J.\ Probab.\
  (2025).

\bibitem{BBD2024} R.~Bauerschmidt, T.~Bodineau, and
  B.~Dagallier, ``Stochastic dynamics and the Polchinski equation:
  An introduction,'' Probab.\ Surv.\ \textbf{21} (2024) 200--290.

\bibitem{BBD2021} R.~Bauerschmidt and T.~Bodineau, ``Log-Sobolev
  inequality for the continuum sine-Gordon model,'' Comm.\ Pure
  Appl.\ Math.\ \textbf{74} (2021) 2064--2113.

\bibitem{CCHS2020} A.~Chandra, I.~Chevyrev, M.~Hairer, and
  H.~Shen, ``Langevin dynamic for the 2D Yang--Mills measure,''
  Publ.\ Math.\ IH\'ES \textbf{136} (2022) 1--147.

\bibitem{CCHS2022} A.~Chandra, I.~Chevyrev, M.~Hairer, and
  H.~Shen, ``Stochastic quantisation of Yang--Mills--Higgs in 3D,''
  Invent.\ Math.\ \textbf{237} (2024) 541--696.

\bibitem{CCHS2024} I.~Chevyrev and H.~Shen, ``Uniqueness of gauge
  covariant renormalisation of stochastic 3D Yang--Mills--Higgs,''
  ARMA (2025).

\bibitem{Balaban1987} T.~Balaban, ``Renormalization group approach
  to lattice gauge field theories,'' Commun.\ Math.\ Phys.\
  \textbf{109} (1987) 249--301.

\bibitem{PW1981} G.~Parisi and Y.-S.~Wu, ``Perturbation theory
  without gauge fixing,'' Sci.\ Sinica \textbf{24} (1981) 483--496.

\bibitem{OS1978} K.~Osterwalder and E.~Seiler, ``Gauge field
  theories on a lattice,'' Ann.\ Phys.\ \textbf{110} (1978)
  440--471.

\end{thebibliography}

\end{document}
